\appendix \label{sec:appendix}
%\section{Bayesian Model Averaging} \label{sec:appendix-bma}   %% Appendix A
%TODO: Describe in statistical literature and how it differs from how it is normallly used with climate models.

\section{Influence of prior on weighted projections}

Figure \ref{fig:appendix-projections-mmm-bma} shows the influence of using a prior distribution over weight vectors in the ensemble performance weighting approach that strongly restricts the considered data. Here, we use a Dirichlet prior with concentration parameter $\alpha=1$ that yields weight vectors with roughly equal weights assigned to every model. This strongly restricts the predictions of the weighted average models that are considered \emph{a priori} as shown in Fig.~\ref{fig:expected-ecs-dirichlet}b in the main text. This, in turn, is reflected in Fig.~\ref{fig:appendix-projections-mmm-bma} where the shown inner 90\% uncertainty range for the projected increase in global mean temperature is nearly identical when using the prior weight vectors as compared to using the posterior weight vectors. Therefore, this extremely narrow range is not due to the data: we get almost the same uncertainty range simply by using the chosen prior weights without any consideration of the likelihood function.

\begin{figure}[H]
    \centering
    %\includegraphics{sections/figures/appendix/projections-bma-mmm-with-prior-1.pdf}
    \includegraphics{sections/figures/figA1.pdf}
    %\caption{Predicted increase in global mean near-surface air temperature with respect to the time period from 1995--2014 until the end of the century, under emission scenario SSP5.85. Solid lines show the multi-model mean (brown) of an ensemble of $N=38$ CMIP6 models and the weighted average of this ensemble (blue) using the mean posterior weights applying the ensemble performance weighting approach with prior Dirichlet($\mathbf{1}$) and with the target distribution from Fig.~\ref{fig:expected-ecs-dirichlet} as likelihood. The green dashed lines show the weighted mean (thick line), using the mean prior weights. Uncertainty bands (for the prior dashed green lines) represent the 0.05 and 0.95 (weighted) quantiles.}
    \caption{Like Fig.~\ref{fig:projections-mmm-bma} in the main text, but when using a Dirichlet prior with $\alpha=1$.}
    \label{fig:appendix-projections-mmm-bma}
\end{figure}

%\subsection{}     %% Appendix A1, A2, etc.

%\footnote{Note that similarly to the case when $\sigma^2$ is sampled in the individual performance weighting approach, when using multiple diagnostics, it may happen that $\sigma_d^2$ for a diagnostic $d$ receives large values so that the observed data for this diagnostic is nearly equally likely under all weighted average models and all  }0

\section{}\label{sec:appendix-B}

Here we show examples of different weighting approaches that yield the same ranking order and only differ with respect to the magnitude of the individual weights. Note that for simplicity, we left out the area-weights here that should be applied for spatial data.

\begin{enumerate}
    \item Assume i.i.d Gaussian errors, i.e. $y(s) = m_i(s) + \epsilon$ with $\epsilon \sim \mathcal{N}(0,\sigma^2)$ and weights for model $i$ are proportional to the likelihood of the data under model $i$
    \begin{itemize}
        \item That is, we assume that each data point (y(s)) comes from a normal distribution with the prediction of model $i$ ($m_i$) as mean and with standard deviation $\sigma$.
        \item The likelihood of the observed data under model $i$ is then given by $P(Y\mid M_i, \sigma) = \prod_{s} {\frac {1}{\sqrt {2\pi \sigma ^{2}}}}\exp {\left(-{\frac {(y(s) - m_i(s)^{2}}{2\sigma ^{2}}}\right)}$, iterating over all data points, e.g. all spatial locations $s$.
        \item The log likelihood of all observed data is then given by: $-{\frac {n}{2}}\log(2\pi \sigma ^{2})-{\frac {1}{2\sigma ^{2}}}\sum _{i=1}^{n}\left(y(s)-m_i(s) \right)^{2}$
        \item turned into a likelihood again: $w_i \propto \exp(c_1) - \exp(\frac{1}{2\sigma^2} \cdot \sum _{i=1}^{n}\left(y(s)-m_i(s) \right)^{2} = \exp(c_1) - \exp(c_2 \cdot \textrm{SSE})$; i.e. $w_i\propto \exp(-c_2\cdot SSE)$ where SSE is the sum of squared errors and $c_1, c_2$ are constants.
    \end{itemize}

\item \citet{sierraInformationtheoreticApproachObtain2025}: $w_i \propto \frac{1}{\hat{\sigma}^2}$ with $\hat{\sigma}^2 = \frac{\sum_{s} (m_i(s) - y(s))^2}{n} = \frac{\textrm{SSE}}{n} = \textrm{MSE}$. Therefore, $w_i \propto n \cdot \frac{1}{\textrm{SSE}}=\frac{1}{\textrm{MSE}}$ with MSE referring to the mean squared error.

\item \citet{brunnerReducedGlobalWarming2020} performance weights for a single diagnostic variable: $w_i \propto \exp(-(\frac{\sqrt{\textrm{MSE}}}{\sigma_D})^2)$ where $\sigma_D$ is a tuned hyperparameter. 

\end{enumerate}


\begin{comment}
\section{Bayesian Model Averaging}\label{sec:appendix-bma}
In the introduction we mentioned that some papers have referred to what we call the ‘Ensemble performance weighting’ approach as ‘Bayesian Model Averaging’, short BMA. However, in the statistical literature, BMA has a different meaning that we believe is worth disentangling and setting in relation to the ensemble performance weighting approach that we have described here. In the classical sense, BMA applies Bayesian inference on the problem of model selection by taking into account the uncertainties about the models themselves instead of assuming one specific model \citep[for a review on BMA see][]{fragosoBayesianModelAveraging2018}. For a concrete statistical model $\mathcal{M}$ with certain parameters, summarized as $\theta_{\mathcal{M}}$, we have the following quantities:

\begin{itemize}
    \item[\textbullet] Posterior distribution: $P(\theta_{\mathcal{M}} \mid Y)$
    \item[\textbullet] Likelihood: $P(Y\mid \theta_{\mathcal{M}})$
    \item[\textbullet] Posterior Predictive Distribution for new data $\tilde{y}$:
        \begin{equation}
            P(\tilde{y}\mid Y, \mathcal{M}) = \int_{\theta_{\mathcal{M}}} P(\hat{y}\mid \theta_{\mathcal{M}}) P(\theta_{\mathcal{M}}\mid Y) d\theta_{\mathcal{M}}
        \end{equation}
\end{itemize}

However, we may not be certain that this model is the best representation for our data. BMA now takes into account this uncertainty about the models themselves when making predictions about new data $\tilde{y}$: 

\begin{equation}
    P(\tilde{y}\mid Y) = \sum_{i=1}^N P(\tilde{y}\mid Y, \mathcal{M}_i) \cdot P(\mathcal{M}_i \mid Y) 
\end{equation}

where the weights, $P(\mathcal{M}_i\mid Y)$, are computed via the marginal likelihood:

\begin{equation}
    P(\mathcal{M}_i\mid Y) \propto P(Y\mid \mathcal{M}_i) = \int_{\theta_{\mathcal{M}}} P(Y\mid \theta_{\mathcal{M},i}) \cdot P(\theta_{\mathcal{M},i}) d\theta_{\mathcal{M}} 
\end{equation}

This, in turn, means that in classical BMA, models are interpreted as representing mutually exclusive hypotheses, including the `correct' (or `most correct') model.  In the long run, after collecting more and more data, all probability mass will be put on the best model. This clashes with the interpretation of an ensemble of climate models which all contribute to our understanding of the climate system and among which we do not believe to find a single `correct' model. Instead, climate models are considered to be inter-dependent imperfect representations of the climate system.
Contrary to the `Ensemble Performance Weighting' approach that we have described, in the classical definition of BMA, the weights are posterior model weights, i.e.~they refer to individual models.

\citet{rafteryUsingBayesianModel2005} extended the statistical definition of BMA to dynamical models, more precisely to ensembles of weather forecast models. They combine the probability densities of the models' forecasts into a mixture distribution via mixture weights that are computed with the Expectation-Maximization (EM) algorithm. The weights are interpreted as the respective forecast to be the best forecast in the ensemble. That is, the 
objective in this 'BMA-approach' is to maximize the joint mixture distribution, however, assuming that each observation originates from one of the component distributions. That is, for every observation, a latent variable is assumed that specifies from which individual model it was generated. The computed weight vector thus summarizes how often each model is the best among all models in the ensemble, in explaining the observed data. 
This contrasts the `Ensemble Performance Weighting' approach which is based on the assumption that a weighted average model generated the observed data and which computes a distribution over weight vectors, instead of estimating a single weight vector. 

There are applications following the classical as well as the approach from \citet{rafteryUsingBayesianModel2005} to BMA in the context climate models \citep[e.g.,][]{zhuFutureProjectionsUncertainty2013,olsonSimpleMethodBayesian2016, fanBayesianPosteriorPredictive2017, zhuIntercomparisonMultimodelEnsembleprocessing2023}.


%\begin{equation}
 %   P(\tilde{y}\mid Y) = \sum_{i=1}^N w_i \cdot P(\tilde{y}\mid \mathcal{M}_i)
%\end{equation}


%Contrary to climate projections, weather forecasts work on very short timescales, thus

\end{comment}


\noappendix       %% use this to mark the end of the appendix section. Otherwise the figures might be numbered incorrectly (e.g. 10 instead of 1).

%% Regarding figures and tables in appendices, the following two options are possible depending on your general handling of figures and tables in the manuscript environment:

%% Option 1: If you sorted all figures and tables into the sections of the text, please also sort the appendix figures and appendix tables into the respective appendix sections.
%% They will be correctly named automatically.

%% Option 2: If you put all figures after the reference list, please insert appendix tables and figures after the normal tables and figures.
%% To rename them correctly to A1, A2, etc., please add the following commands in front of them:

\appendixfigures  %% needs to be added in front of appendix figures



\appendixtables   %% needs to be added in front of appendix tables

%% Please add \clearpage between each table and/or figure. Further guidelines on figures and tables can be found below.

